% =============================================================
% layout.tex - Seitenränder, Abstände, Kopf-/Fußzeilen, Verzeichnisse
% Ersetzt: Seitenrand-settings.tex + Absatz_Abstand-settings.tex +
%          Kopf_und_Fusszeile_misc.tex + Tocstyle.tex +
%          Abkuerzungsverzeichnis-settings.tex
% =============================================================

% =============================================================
% 1. SEITENRÄNDER
% =============================================================

\usepackage{geometry}

\newif\ifdigitalversion
\digitalversiontrue  % Für PDF/Digital (auf false setzen für Druck)

\ifdigitalversion
  \geometry{
    a4paper, twoside,
    left=30mm, right=30mm,
    top=25mm, bottom=30mm,
    headheight=15pt, headsep=10mm, footskip=12mm
  }
\else
  \geometry{
    a4paper, twoside,
    left=35mm, right=25mm,
    top=25mm, bottom=30mm,
    headheight=15pt, headsep=10mm, footskip=12mm,
    bindingoffset=5mm
  }
\fi

\raggedbottom

% =============================================================
% 2. ABSATZABSTÄNDE & LISTEN
% =============================================================

\setlength{\emergencystretch}{3em}

% Listen-Konfiguration
\setlist{parsep=0pt, itemsep=3pt, topsep=6pt, partopsep=0pt}
\setlist[itemize]{label=\textcolor{primary}{\textbullet}, leftmargin=1.5em}
\setlist[enumerate]{label=\textcolor{secondary}{\arabic*.}, leftmargin=2em}
\setlist[description]{font=\bfseries\color{headingcolor}, leftmargin=2em}

% Utility-Befehle für Absatzabstände
\newcommand{\tightpars}{\setlength{\parskip}{3pt plus 1pt minus 1pt}}
\newcommand{\loosepars}{\setlength{\parskip}{9pt plus 3pt minus 2pt}}
\newcommand{\resetpars}{\KOMAoptions{parskip=half}}

% =============================================================
% 3. KOPF- UND FUSSZEILEN
% =============================================================

\usepackage[automark, headsepline]{scrlayer-scrpage}
\pagestyle{scrheadings}
\clearscrheadings
\clearscrplain

% Gerade Seiten (links)
\lehead{\rightmark}
\rehead{}
% Ungerade Seiten (rechts)
\lohead{}
\rohead{\leftmark}
% Seitenzahlen außen
\lefoot{\pagemark}
\refoot{}
\lofoot{}
\rofoot{\pagemark}
% Plain-Stil (Kapitelanfänge)
\lefoot[\pagemark]{\pagemark}
\rofoot[\pagemark]{\pagemark}

\automark[section]{chapter}

% =============================================================
% 4. INHALTSVERZEICHNIS (TOC)
% =============================================================

% Schriftarten für TOC-Elemente
\addtokomafont{disposition}{\bfseries}
\addtokomafont{chapter}{\LARGE\bfseries\color{headingcolor}}
\addtokomafont{section}{\Large\bfseries\color{headingcolor}}
\addtokomafont{subsection}{\large\bfseries\color{headingcolor}}
\addtokomafont{subsubsection}{\normalsize\bfseries\color{headingcolor}}
\addtokomafont{paragraph}{\normalsize\bfseries\color{headingcolor}}

% Titelseite
\addtokomafont{title}{\huge\bfseries\color{headingcolor}}
\addtokomafont{author}{\Large\color{secondary}}
\addtokomafont{date}{\large\color{secondary}}

% Kopf/Fuß, Beschriftungen
\addtokomafont{pageheadfoot}{\small\color{secondary}}
\addtokomafont{pagenumber}{\small\bfseries}
\addtokomafont{caption}{\small\color{secondary}}
\addtokomafont{captionlabel}{\small\bfseries\color{headingcolor}}
\addtokomafont{footnote}{\footnotesize}
\addtokomafont{chapterentry}{\bfseries\color{headingcolor}}

% TOC-Konfiguration
\KOMAoptions{toc=chapterentrywithdots}

% Deutsche Bezeichnungen
\renewcommand{\contentsname}{Inhaltsverzeichnis}
\renewcommand{\listfigurename}{Abbildungsverzeichnis}
\renewcommand{\listtablename}{Tabellenverzeichnis}
\renewcommand{\figurename}{Abbildung}
\renewcommand{\tablename}{Tabelle}
\renewcommand{\chaptername}{Kapitel}
\renewcommand{\appendixname}{Anhang}

% =============================================================
% 5. ABKÜRZUNGSVERZEICHNIS (GLOSSARIES)
% =============================================================

% Deutsche Bezeichnungen
\renewcommand{\glossaryname}{Glossar}
\renewcommand{\acronymname}{Abkürzungsverzeichnis}

% Stil
\setglossarystyle{long3col}
\setacronymstyle{long-short}
\renewcommand*{\glstextformat}[1]{\textcolor{primary}{#1}}
\setlength{\glsdescwidth}{0.6\textwidth}
\setlength{\glspagelistwidth}{0.1\textwidth}

% Glossar-Befehle
\renewcommand*{\glsnamefont}[1]{\textbf{\textcolor{headingcolor}{#1}}}
\setglossarysection{section}

% Hilfs-Befehle
\newcommand{\abk}[2]{\newacronym{#1}{#1}{#2}}
\newcommand{\begriff}[3][]{\newglossaryentry{#2}{name={#2}, description={#3}, #1}}

% Zusätzliche Glossar-Arten
\newglossary[slg]{symbols}{syi}{syg}{Symbolverzeichnis}

% Glossar aktivieren
\makeglossaries
